\documentclass[11pt]{article}

\usepackage[a4paper,margin=1in]{geometry}
\usepackage{amsmath,amssymb,amsthm}
\usepackage{booktabs}
\usepackage{hyperref}

\newcommand{\R}{\mathbb{R}}
\newcommand{\I}{\mathbb{I}}
\newcommand{\E}{\mathbb{E}}
\newcommand{\dd}{\,\mathrm{d}}
\newcommand{\given}{\,\mid\,}
\newcommand{\npcc}{\mathrm{NPCC}}
\newcommand{\logit}{\operatorname{logit}}

\title{Mathematical Details of the Neural Pair-Copula Construction}
\author{}
\date{}

\begin{document}
\maketitle

\section{Setup}

Let
\[
  (U_i,V_i,X_i)_{i=1}^n,
  \qquad
  U_i,V_i \in (0,1), \quad X_i \in \R^p,
\]
denote observations on the copula scale, possibly with covariates.
The goal is to estimate the conditional bivariate copula density
\[
  c(u,v \given x),
  \qquad (u,v) \in (0,1)^2.
\]
When no covariates are supplied, the implementation uses a constant
one-dimensional covariate, so the unconditional case is represented as
a special case of the conditional formulation.

\section{Rosenblatt Factorization}

On the copula scale, the conditional margins of $U$ and $V$ are uniform.
Hence a conditional bivariate copula density may be represented through a
Rosenblatt factorization,
\[
  c(u,v \given x)
  =
  f_{V \given U,X}(v \given u,x).
\]
The reverse ordering gives the analogous identity
\[
  c(u,v \given x)
  =
  f_{U \given V,X}(u \given v,x).
\]
Each factor is a univariate conditional density. The estimator therefore
reduces conditional copula density estimation to two conditional
univariate predictive-distribution problems.

\section{Neural Pair-Copula Estimator}

The model fits two TabPFN-backed conditional distribution estimators:
\[
  \widehat f_{V \given U,X}(v \given u,x),
  \qquad
  \widehat f_{U \given V,X}(u \given v,x).
\]
For the first direction, the conditioning features are
\[
  W^{(V)} = (U,X),
\]
and the response is $V$. For the reverse direction,
\[
  W^{(U)} = (V,X),
\]
and the response is $U$.

The fitted Neural Pair-Copula Construction (NPCC) density is the
symmetric average
\[
  \widehat c_{\npcc}(u,v \given x)
  =
  \frac{1}{2}\widehat f_{V \given U,X}(v \given u,x)
  +
  \frac{1}{2}\widehat f_{U \given V,X}(u \given v,x).
  \label{eq:symmetric-density}
\]
This averaging reduces sensitivity to the Rosenblatt ordering. Without
the optional projection in Section~\ref{sec:sinkhorn}, the average is
not guaranteed to satisfy the two uniform marginal constraints exactly.

\section{Support Transforms}

The inner univariate model can fit a transformed response
\[
  Z = T(Y),
  \qquad Y \in (0,1),
\]
where $Y$ is either $U$ or $V$. Densities are converted back to the
copula scale by the change-of-variables formula
\[
  f_Y(y \given w)
  =
  f_Z(T(y) \given w)\,\left|T'(y)\right|.
  \label{eq:change-of-variables}
\]
The implementation supports the transformations in
Table~\ref{tab:transforms}.

\begin{table}[h]
  \centering
  \begin{tabular}{lll}
    \toprule
    Name & Transform $T(y)$ & Jacobian $\left|T'(y)\right|$ \\
    \midrule
    identity & $y$ & $1$ \\
    logit & $\logit(y)=\log\{y/(1-y)\}$ & $\{y(1-y)\}^{-1}$ \\
    probit & $\Phi^{-1}(y)$ & $\phi(\Phi^{-1}(y))^{-1}$ \\
    \bottomrule
  \end{tabular}
  \caption{Response transformations supported by the inner conditional
  distribution model. Here $\Phi$ and $\phi$ are the standard normal CDF
  and PDF.}
  \label{tab:transforms}
\end{table}

The logit transform is the default. Inputs are clipped away from the
boundary of $(0,1)$ by a small numerical tolerance before applying the
transform.

\section{TabPFN Predictive Distribution Recovery}

Both inner estimators train a TabPFN regressor on pairs
$(W_i,T(Y_i))$. They differ only in how the fitted conditional
predictive distribution is read from TabPFN.

\subsection{Criterion Method}

The default method uses TabPFN's native distribution head. For a query
feature matrix $W$, TabPFN returns logits $\ell(W)$ over a binned
predictive distribution and a criterion object that evaluates the
corresponding conditional density, distribution function, and quantile
function. On the transformed scale,
\[
  \widehat f_Z(z \given w)
  =
  \operatorname{criterion.pdf}(\ell(w), z),
\]
\[
  \widehat F_Z(z \given w)
  =
  \operatorname{criterion.cdf}(\ell(w), z),
\]
and
\[
  \widehat F_Z^{-1}(\alpha \given w)
  =
  \operatorname{criterion.icdf}(\ell(w), \alpha).
\]
The density on the copula scale follows from
Equation~\eqref{eq:change-of-variables}; CDFs and quantiles are mapped
through the monotone transform without a Jacobian correction.

\subsection{Quantile Method}

The alternative method queries TabPFN for a conditional quantile table
\[
  \widehat Q_Z(\alpha_k \given w),
  \qquad
  \alpha_k \in [\alpha_{\min},\alpha_{\max}],
  \quad k=1,\dots,K.
\]
For each query row, the predicted quantiles are sorted to enforce
monotonicity. The conditional CDF is recovered by interpolation in the
quantile table:
\[
  \widehat F_Z(z \given w)
  \approx
  \alpha
  \quad\text{such that}\quad
  \widehat Q_Z(\alpha \given w) = z.
\]
The conditional density is recovered by differentiating the quantile
function,
\[
  \widehat f_Z(z \given w)
  \approx
  \left[
    \frac{\partial}{\partial \alpha}
    \widehat Q_Z(\alpha \given w)
  \right]^{-1}_{\alpha=\widehat F_Z(z \given w)}.
  \label{eq:quantile-density}
\]
In the implementation, the derivative is computed by finite differences,
floored by a positive constant for numerical stability, and then linearly
interpolated.

\section{Conditional Distribution Functions}

The fitted model exposes two conditional distribution functions, also
called $h$-functions:
\[
  \widehat h_1(u,v \given x)
  =
  \widehat F_{V \given U,X}(v \given u,x),
  \label{eq:hfunc1}
\]
\[
  \widehat h_2(u,v \given x)
  =
  \widehat F_{U \given V,X}(u \given v,x).
  \label{eq:hfunc2}
\]
The numbering follows the pyvinecopulib convention: $h_1$ conditions on
the first argument and $h_2$ conditions on the second argument.

The joint conditional copula CDF is approximated by averaging the two
integrated Rosenblatt directions,
\[
  \widehat C(u,v \given x)
  =
  \frac{1}{2}
  \left[
    \int_0^u
    \widehat F_{V \given U,X}(v \given s,x)\dd s
    +
    \int_0^v
    \widehat F_{U \given V,X}(u \given t,x)\dd t
  \right].
  \label{eq:joint-cdf}
\]
The implementation evaluates these integrals by the trapezoidal rule.

\section{Optional Sinkhorn Projection}
\label{sec:sinkhorn}

The symmetric density in Equation~\eqref{eq:symmetric-density} is
nonnegative but does not necessarily satisfy the copula marginal
constraints exactly:
\[
  \int_0^1 c(u,v \given x)\dd v = 1,
  \qquad
  \int_0^1 c(u,v \given x)\dd u = 1.
\]
An optional iterative proportional fitting, or Sinkhorn, projection is
available to enforce these constraints approximately on a grid.

Let $u_1,\dots,u_m$ and $v_1,\dots,v_q$ be projection grids with
trapezoidal weights $a_i$ and $b_j$. For a fixed covariate value $x$,
write the raw density matrix as
\[
  C_{ij} = \widehat c_{\npcc}(u_i,v_j \given x).
\]
The projection seeks positive row and column scalings $r_i$ and $s_j$
such that
\[
  \widetilde C_{ij} = r_i C_{ij} s_j
\]
has approximately uniform margins under the quadrature rule:
\[
  \sum_{j=1}^q \widetilde C_{ij} b_j \approx 1,
  \qquad
  \sum_{i=1}^m \widetilde C_{ij} a_i \approx 1.
\]
The implementation alternates the updates
\[
  r_i
  \leftarrow
  \left(\sum_{j=1}^q C_{ij}s_j b_j\right)^{-1},
  \qquad
  s_j
  \leftarrow
  \left(\sum_{i=1}^m C_{ij}r_i a_i\right)^{-1},
\]
computed in the log domain for numerical stability. For pointwise
density queries, the grid scalings are linearly interpolated back to the
requested $(u,v)$ locations.

\section{Kendall's Tau}

For a fixed covariate value $x$, Kendall's tau is estimated by a
simulation recipe based on the inverse Rosenblatt transform. Draw a
deterministic two-dimensional generalized Halton sample
\[
  (A_i,B_i)_{i=1}^N \subset (0,1)^2.
\]
Set
\[
  U_i^\star = A_i,
  \qquad
  V_i^\star =
  \widehat F^{-1}_{V \given U,X}(B_i \given A_i,x).
\]
Then $(U_i^\star,V_i^\star)$ is an approximate sample from the fitted
copula, and Kendall's tau is computed as the rank Kendall correlation of
these pairs.

\section{Simulation Study Scenarios}

The simulation utilities generate data from known pyvinecopulib
bivariate copula families. The implemented one-parameter families are
Clayton, Gumbel, Frank, Gaussian, and Joe. Each family is parameterized
through Kendall's tau.

For conditional scenarios, a scalar covariate $x \in [0.01,0.99]$
controls Kendall's tau through one of the following regimes:
\[
  \tau_{\mathrm{linear}}(x)
  =
  \tau_{\min}+(\tau_{\max}-\tau_{\min})x,
\]
\[
  \tau_{\mathrm{constant}}(x)=0.5,
\]
\[
  \tau_{\mathrm{sin}}(x)
  =
  \operatorname{clip}\{0.5+0.4\sin(2\pi x)\},
\]
\[
  \tau_{\mathrm{quadratic}}(x)
  =
  \tau_{\min}+(\tau_{\max}-\tau_{\min})(2x-1)^2,
\]
where the implementation clips the tau values to
$[\tau_{\min},\tau_{\max}]=[0.05,0.90]$. The unconditional scenarios use
fixed Kendall's tau values $0.25$, $0.50$, and $0.75$.

\section{Implementation Correspondence}

\begin{table}[h]
  \centering
  \begin{tabular}{ll}
    \toprule
    Mathematical object & Implementation \\
    \midrule
    Symmetric estimator in Equation~\eqref{eq:symmetric-density}
      & \texttt{PFNRBicop} \\
    Criterion predictive distribution
      & \texttt{TabPFNCriterionDistribution1D} \\
    Quantile predictive distribution
      & \texttt{TabPFNQuantileDistribution1D} \\
    Shared support transforms
      & \texttt{TabPFNDistribution1D} \\
    Sinkhorn projection
      & \texttt{sinkhorn\_iters} \\
    \bottomrule
  \end{tabular}
\end{table}

\end{document}
